% Compiler avec XeLaTeX ou LuaLaTeX
\documentclass[11pt,a4paper]{article}

\usepackage{amsmath,amssymb}
\usepackage{fontspec}
\usepackage{unicode-math}
\usepackage{polyglossia}
\setmainlanguage{french}
\usepackage{xcolor}
\usepackage{fancyhdr}
\usepackage[margin=2cm,headheight=16pt,headsep=22pt]{geometry}
\usepackage{enumitem}

% Police principale et police mathématique : Libertinus
\IfFontExistsTF{Libertinus Serif}
  {\setmainfont{Libertinus Serif}}
  {\IfFontExistsTF{Latin Modern Roman}{\setmainfont{Latin Modern Roman}}{\setmainfont{lmroman10-regular.otf}}}

\IfFontExistsTF{Libertinus Math}
  {\setmathfont{Libertinus Math}}
  {\IfFontExistsTF{Latin Modern Math}{\setmathfont{Latin Modern Math}}{\setmathfont{latinmodern-math.otf}}}

\linespread{1.07}
\setlength{\parindent}{0pt}
\AtBeginDocument{\setlength{\parindent}{0pt}}
\setlength{\parskip}{0.35em}

% En-tête
\pagestyle{fancy}
\fancyhf{}
\lhead{NeoHarness AI}
\rhead{Lycée -- Mathématiques}
% Pagination : « Page X / Y » centrée en pied de page
\usepackage{lastpage}
\cfoot{Page \thepage\ / \pageref{LastPage}}
\renewcommand{\headrulewidth}{0.5pt}
% Trait de séparation éloigné du texte de l'en-tête (modifier 5pt pour ajuster)
\renewcommand{\headrule}{\vskip 5pt\hrule width\headwidth height\headrulewidth \vskip-\headrulewidth}

% Configuration des listes
\setlist[enumerate,1]{itemsep=0.45em, topsep=0.3em, label=\textbf{\arabic*.}}
\setlist[enumerate,2]{itemsep=0.25em, topsep=0.2em, label=\textbf{\alph*.}}

% Commande exercice
\newcommand{\exo}[2]{%
  \vspace{1.5em}%
  \noindent{\large\textbf{Exercice #1. -- #2}}%
  \par\vspace{0.35em}%
  \hrule height 0.5pt%
  \vspace{0.7em}%
}

\newcommand{\partie}[1]{%
  \vspace{0.9em}%
  \noindent\textbf{#1}%
  \par\vspace{0.3em}%
}

\begin{document}

\begin{center}
  {\LARGE \textbf{Exercice : Produit scalaire dans le plan et espace}}\\[0.5em]
  {\large Lycée -- Mathématiques}
\end{center}

Le produit scalaire est un outil fondamental de la géométrie vectorielle au lycée, dont l'étude s'articule rigoureusement entre le plan (en Première spécialité mathématiques) et l'espace (en Terminale). Il permet de faire le pont entre l'algèbre (coordonnées, repères orthonormés) et la métrique (longueurs, orthogonalité, angles).

\vspace{0.35cm}\noindent\textbf{\large 1. Le produit scalaire dans le plan (Première Spécialité)}\par\vspace{0.15cm}
Dans le plan euclidien orienté muni d'un repère orthonormé, le produit scalaire de deux vecteurs $\vec{u}$ et $\vec{v}$, noté $\vec{u} \cdot \vec{v}$, se définit de plusieurs manières complémentaires :
\begin{enumerate}
  \item \textbf{Définition géométrique (projeté orthogonal)} : Si $\vec{u} = \overrightarrow{AB}$ et $\vec{v} = \overrightarrow{AC}$, en introduisant le projeté orthogonal $H$ de $C$ sur la droite $(AB)$, on a $\overrightarrow{AB} \cdot \overrightarrow{AC} = AB \times AH$ (avec un signe positif si $H$ est dans le sens de $\vec{AB}$, négatif sinon).
  \item \textbf{Définition trigonométrique} : $\vec{u} \cdot \vec{v} = \|\vec{u}\| \times \|\vec{v}\| \times \cos(\vec{u}, \vec{v})$.
  \item \textbf{Expression analytique} : Si $\vec{u}(x ; y)$ et $\vec{v}(x' ; y')$, alors $\vec{u} \cdot \vec{v} = xx' + yy'$.
  \item \textbf{Orthogonalité} : Deux vecteurs sont orthogonaux si et seulement si leur produit scalaire est nul.
\end{enumerate}

\vspace{0.35cm}\noindent\textbf{\large 2. Extension au produit scalaire dans l'espace (Terminale)}\par\vspace{0.15cm}
En classe de Terminale, le cadre s'étend à l'espace tridimensionnel muni d'un repère orthonormal $(O ; \vec{i}, \vec{j}, \vec{k})$.
\begin{enumerate}
  \item \textbf{Expression analytique} : Pour $\vec{u}(x ; y ; z)$ et $\vec{v}(x' ; y' ; z')$, le produit scalaire s'exprime par la formule générale :
\end{enumerate}
$$\vec{u} \cdot \vec{v} = xx' + yy' + zz'$$
\begin{enumerate}
  \item \textbf{Propriétés algébriques} : Le produit scalaire demeure bilinéaire et symétrique ($\vec{u} \cdot \vec{v} = \vec{v} \cdot \vec{u}$).
  \item \textbf{Vecteur normal et équation de plan} : Un vecteur non nul $\vec{n}$ est normal à un plan $\mathcal{P}$ s'il est orthogonal à deux vecteurs directeurs non colinéaires de ce plan. Cette propriété permet de déterminer directement l'équation cartésienne d'un plan de la forme $ax + by + cz + d = 0$, où $(a ; b ; c)$ sont les coordonnées de $\vec{n}$.
\end{enumerate}

\end{document}
